\begin{corollary}[L2 load split]\label{cor:L2-load-split}
The \(L2\) positive clause is the finite-fibre certificate construction; the
strictness clause is the separate transfer obstruction of \Cref{thm:finite-state-strictness}.
\end{corollary}

\begin{proof}
The positive clause is \Cref{prop:zeckendorf-cert-univ}.  The obstruction is the
failure of the finite-state budget, which is not a fixed-resolution moment
claim by
\Cref{prop:L2-novelty-finite-state-separation,thm:L2-finite-window-transfer-obstruction}.
\end{proof}

\begin{proposition}[Terminology bridge for the \texorpdfstring{\(L2\)}{L2} clause]
\label{thm:L2-terminology-bridge}
In this paper, \(L2\)-certificate universality means finite-fibre certificate
universality in the sense of \Cref{def:cert-univ}.  A finite-state transfer
formula is a strictly additional hypothesis on that certificate data.  Hence the
\(L2\) strictness statements are exactly cover-relative failures of this extra
finite-state transfer hypothesis, not failures of finite-fibre certificate
universality itself.
\end{proposition}

\begin{proof}
The first sentence is the definition of \(L2\)-certificate universality in
\Cref{def:cert-univ}: it asks for effective finite-resolution data whose
fibre-size multisets supply the fixed finite certificates of
\Cref{prop:classical-finite-atomic-skeleton}.  The same definition then states
finite-state transfer separately, by requiring for each fixed \(q\) a single
matrix expression \(S_q^{(m)}=u_q^{\mathsf T}A_q^m v_q\) valid for all
resolutions \(m\).  \Cref{prop:L2-novelty-finite-state-separation} proves that
this uniform-in-\(m\) matrix condition is not contained in the fixed-resolution
certificate data, and \Cref{cor:L2-load-split} identifies the positive clause
with \Cref{prop:zeckendorf-cert-univ} while assigning strictness to the transfer
budget obstruction.  Finally,
\Cref{thm:L2-finite-fiber-no-transfer-content} proves the needed obstruction
inside the computable finite-fibre cover category over the same normal forms.
Thus the terminology has the asserted unique reading.
\end{proof}

\begin{corollary}[Subordinate role of L2 finite-fibre facts]
\label{cor:L2-subordinate-role}
The finite-fibre facts do not prove the Richardson obstruction, the one-orbit
multiplication obstruction, or the Zeckendorf \(L1\) protocol.  They only attach
finite certificates after a normalizer has already been chosen.  In particular,
no result labelled in \Cref{sec:consequences} is a premise for any theorem in
\Cref{sec:L0,sec:L1,sec:strictness}.
\end{corollary}

\begin{proof}
Every result in this section starts from finite fibres of a given normalizing
map.  The \(L1\), Richardson, and one-orbit arguments are proved in
\Cref{sec:L1,sec:strictness} and do not use these certificates.
\end{proof}

\begin{theorem}[Cover moments do not determine intrinsic fold transfer]
\label{thm:L2-cover-intrinsic-nonimplication}
Let
\[
  d_m(x)=|\Fold_m^{-1}(x)|,
  \qquad S_2(m)=\sum_{x\in X_m}d_m(x)^2
\]
be the intrinsic second collision moment of the unmodified Zeckendorf fold.  Let
\(d_m^\sharp(x)\) be the fibre-size function of a computable finite-fibre cover
over the same normal-form set \(X_m\), and let
\[
  S_2^\sharp(m)=\sum_{x\in X_m}d_m^\sharp(x)^2 .
\]
The assertion that \((S_2^\sharp(m))\) has, or fails to have, a finite-state
transfer formula is a statement about the cover-size data \(d_m^\sharp\).  It
does not imply the corresponding assertion for the intrinsic sequence
\((S_2(m))\) unless an additional theorem identifies the two sequences, or gives
a transfer-preserving comparison between them, for all \(m\).  No such
identification or comparison is part of the \(L2\) strictness proof in this
article.
\end{theorem}

\begin{proof}
Both moments are finite sums over the same finite set \(X_m\), but their summands
come from different data.  The intrinsic summand \(d_m(x)\) is defined by the
original map \(\Fold_m\colon\Omega_m\to X_m\).  A computable finite-fibre cover,
by contrast, is specified by a separate computable positive integer function
\(d_m^\sharp(x)\) over \(X_m\); by \Cref{def:computable-cover} and
\Cref{thm:finite-fiber-cover-realization}, this function may be chosen without
changing \(X_m\), \(\Fold_m\), or any readout on normal forms.  Thus the cover
moment can be altered while the intrinsic fold moment remains the moment of the
unchanged map.

Finite-state transfer is a property of a particular sequence in \(m\): by
\Cref{def:cert-univ}, it asks for a fixed representation
\(u^{\mathsf T}A^m v\) for that sequence.  A representation, recurrence, or
growth obstruction for \((S_2^\sharp(m))\) therefore concerns only the sequence
computed from \(d_m^\sharp\).  It can be transported to \((S_2(m))\) only after a
separate equality or transfer-preserving comparison between the two sequences has
been proved.  The one-spike construction of
\Cref{thm:L2-finite-fiber-no-transfer-content} proves the opposite structural
fact needed here: it deliberately changes one cover fibre over an otherwise
fixed normal-form system.  Hence the cover-relative finite-state obstruction has
no intrinsic Fold\(_m\) nontransfer consequence in the present theorem chain; the
positive intrinsic transfer assertion is the separate theorem
\Cref{thm:intrinsic-fold-moment-transfer}.
\end{proof}

\begin{proposition}[Cover-category invariance of the finite-state obstruction]
\label{prop:cover-category-invariance-finite-state-obstruction}
\label{thm:cover-relative-transfer-is-not-intrinsic}
Fix the unmodified Zeckendorf normal-form system \((X_m,\Fold_m)_{m\ge1}\).  The
assignment
\[
  d_m^\sharp\longmapsto S_2^\sharp(m)=\sum_{x\in X_m}d_m^\sharp(x)^2
\]
from computable finite-fibre cover-size functions to second covered moments is
not determined by \((X_m,\Fold_m)\): over this same intrinsic system there is a
singleton computable cover whose second moment has fixed finite-state transfer,
and there is a factorial one-spike computable cover whose second moment has no
fixed finite-state transfer.  Hence finite-state transfer or nontransfer of
covered moments is not an intrinsic invariant of the fold map; it is an
additional property of the chosen cover-size data.
\end{proposition}

\begin{proof}
For the singleton cover take \(d_m^\sharp(x)=1\) for every \(x\in X_m\).  Then
\(S_2^\sharp(m)=|X_m|=F_{m+2}\), and the Fibonacci recurrence gives a fixed
finite-state transfer formula.  For the second cover, choose one normal form in
each \(X_m\), give it cover size \(m!+1\), and give all other normal forms cover
size \(1\).  The construction is computable and changes no normal form or
readout, so it is a cover of the same intrinsic system.  Its second moment is
\((m!+1)^2+F_{m+2}-1\), which grows faster than every exponential.  By
\Cref{prop:finite-state-coefficient-budget}, every fixed finite-state transfer
sequence is exponentially bounded; hence this second covered moment has no fixed
finite-state transfer.

Both covers lie over the identical data \((X_m,\Fold_m)\).  Since they have
opposite transfer behaviour, that behaviour cannot be a function of the
intrinsic fold map alone.  It is therefore exactly a cover-relative property of
the auxiliary multiplicities \(d_m^\sharp\).
\end{proof}
\begin{theorem}[Cover finite-fibre certificates do not imply transfer]
\label{thm:L2-finite-fiber-no-transfer-content}
\label{thm:cover-relative-finite-fibre-no-transfer}
In the Zeckendorf computable finite-fibre cover category, and only in that
cover-relative category, finite-fibre certificate data alone impose no
recurrence, rational generating series, exponential growth bound, or fixed
finite-matrix transfer formula on the cover moments \((S_q^\sharp(m))_{m\ge1}\).
More precisely, for every computable sequence \(a_m\ge1\) there is a computable
finite-fibre cover preserving the normal forms and readouts whose second moment is
\[
  S_2^\sharp(m)=(a_m+1)^2+F_{m+2}-1 .
\]
Taking \(a_m=m!\) gives a computable finite-fibre cover certificate system with
no finite-state transfer formula for \(S_2^\sharp(m)\).  No intrinsic failure
theorem for the unmodified \(\Fold_m\) fibre multiplicities is asserted.
\end{theorem}

\begin{proof}
At resolution \(m\), choose one normal form and assign its cover fibre size
\(a_m+1\), while assigning cover fibre size \(1\) to every other normal form.
This finite cover does not change the normal-form set or any readout on normal
forms, and it is computable because \((a_m)\) is computable.  The displayed
formula follows by adding the square of the enlarged fibre and the squares of
the remaining singleton cover fibres.  If \(a_m=m!\), then
\((S_2^\sharp(m))\) grows faster than
\(C\rho^m\) for every \(C,
\rho\), contradicting the exponential budget in
\Cref{prop:finite-state-coefficient-budget}.  Hence no fixed finite-state
transfer formula exists.
\end{proof}

\begin{theorem}[Two-cover skeleton for the intrinsic boundary]
\label{thm:L2-two-cover-intrinsic-boundary}
Fix the unmodified Zeckendorf data \((X_m,\Fold_m)\).  There are two computable
finite-fibre covers of this same data, both preserving all normal forms and all
\(L1\) readouts, with opposite finite-state-transfer behaviour for their second
collision moments:
\begin{enumerate}
\item the singleton cover has
  \(S_{2,0}^\sharp(m)=\lvert X_m\rvert=F_{m+2}\), hence satisfies a fixed
  finite-state transfer formula; and
\item the factorial one-spike cover has
  \(S_{2,!}^\sharp(m)=(m!+1)^2+F_{m+2}-1\), hence has no fixed finite-state
  transfer formula.
\end{enumerate}
Consequently finite-state transfer or failure for cover moments is not an
intrinsic invariant of the original \(\Fold_m\) fibres.  Any theorem about
recurrence, nonrecurrence, rational generating series, or finite-state transfer
for the intrinsic \(\Fold_m\) multiplicity sequence requires an additional bridge
identifying that intrinsic sequence with the chosen cover sequence; no such
bridge is part of the present article.
\end{theorem}

\begin{proof}
For the first cover, attach one certificate point over every normal form.  The
second collision moment is then the number of normal forms, namely
\(\lvert X_m\rvert=F_{m+2}\) by \Cref{prop:finite-zeckendorf-residue-interface}.  The Fibonacci
sequence satisfies the fixed two-state recurrence
\[
  \binom{F_{m+2}}{F_{m+3}}
  =
  \begin{pmatrix}0&1\\ 1&1\end{pmatrix}^{m}
  \binom{F_2}{F_3},
\]
after the harmless shift of the initial vector, so this moment sequence has a
finite-state transfer formula.

For the second cover, apply \Cref{thm:L2-finite-fiber-no-transfer-content} with
\(a_m=m!\).  It changes only one auxiliary certificate fibre over a fixed normal
form, preserves all readouts by construction, and has displayed moment
\((m!+1)^2+F_{m+2}-1\).  This sequence grows faster than every exponential, while
\Cref{prop:finite-state-coefficient-budget} gives an exponential bound for every
fixed finite-state transfer sequence; hence the factorial one-spike moment has
no finite-state transfer formula.

Both covers project to the identical unmodified data \((X_m,\Fold_m)\).  Since a
property of the auxiliary cover moment takes different truth values on two
covers with the same intrinsic fold map, that property is not determined by the
intrinsic fold multiplicities.  A conclusion for the intrinsic sequence would
therefore require an added equality or transfer-preserving comparison between
that intrinsic sequence and the chosen cover moment, exactly the bridge excluded
in \Cref{thm:L2-cover-intrinsic-nonimplication}.
\end{proof}

\begin{proposition}[Finite-fibre certificate skeleton and transfer pairing]
\label{thm:L2-fiber-multiset-factorization}
\label{thm:L2-finite-fiber-transfer-pairing}
For every fixed resolution \(m\), the \(L2\) certificates used in this paper
factor through the fibre-size multiset \(\mathcal D_m\).  If the same certificate
system also has a fixed finite-state transfer formula for \(S_q^{(m)}\), then
\((S_q^{(m)})_m\) satisfies a constant-coefficient recurrence and a rational
generating series.  Conversely, there are computable one-spike covers preserving
the same normal forms and readouts whose fixed-resolution finite-fibre
certificates exist at every \(m\) but whose second collision moments admit no
fixed finite-state transfer formula.
\end{proposition}

\begin{proof}
The fixed-resolution factorization is exactly
\Cref{prop:finite-fiber-certificate-taxonomy}: the finite-fibre certificates used
here are functions of \(\mathcal D_m\), hence are unchanged by replacing a datum
with another datum having the same fibre-size multiset.  If a fixed finite-state transfer formula
exists, \Cref{def:fixed-finite-state-transfer} gives
\(S_q^{(m)}=u_q^{\mathsf T}A_q^m v_q\) with one finite matrix model independent of
\(m\).  Applying \Cref{prop:finite-state-coefficient-budget} yields the stated
constant-coefficient recurrence and rational generating series.

For the converse, use \Cref{thm:L2-finite-fiber-no-transfer-content} with
\(a_m=m!\).  The construction changes only one auxiliary cover fibre over a fixed
normal form, so it preserves the normal-form set and all readouts, while its
second moment is \((m!+1)^2+F_{m+2}-1\).  This sequence grows faster than every
exponential sequence, contradicting the exponential consequence of any fixed
finite-state transfer formula in \Cref{prop:finite-state-coefficient-budget}.
Thus finite-fibre certificates and fixed finite-state transfer are paired only
when the additional transfer structure is explicitly supplied.
\end{proof}

\begin{corollary}[Cover finite-fibre certificates are independent of transfer]
\label{cor:hankel-psd-independent}
In the Zeckendorf cover category, the fixed-resolution finite-fibre certificates
used here can hold at every resolution while finite-state transfer fails.
\end{corollary}

\begin{proof}
They hold for every finite-fibre system by
\Cref{prop:classical-finite-atomic-skeleton}; finite-state transfer fails for the
factorial one-spike cover in \Cref{thm:L2-finite-fiber-no-transfer-content}.
\end{proof}


